% ============================================================
% paperflow article template — single-column SCI journal style
% All double-at placeholders are filled by the paperflow pipeline
% via plain Python string replacement (contract: see README.md).
% Do NOT rename them.
% Compiles with: tectonic template_filled.tex
% Only TeX Live core packages are used (no shell-escape, no minted).
% ============================================================
\documentclass[11pt,a4paper]{article}

% ---------- page layout ----------
\usepackage[margin=2.5cm]{geometry}

% ---------- math (load amsmath BEFORE newtxmath) ----------
\usepackage{amsmath}

% ---------- Times-like text and math fonts ----------
\usepackage{newtxtext}
\usepackage{newtxmath}

% ---------- graphics and tables ----------
\usepackage{graphicx}
\usepackage{booktabs}

% ---------- captions ----------
\usepackage[labelfont=bf,font=small]{caption}

% ---------- citations: author-year, (Author, 2024; Other, 2025) ----------
\usepackage[round,semicolon,authoryear]{natbib}

% ---------- line spacing ----------
\usepackage{setspace}
\setstretch{1.15}

% ---------- hyperref MUST be loaded last ----------
\usepackage[hidelinks]{hyperref}

% ---------- title block ----------
\title{Multimodal BSE-EDS Fusion for Micro-phase Identification and Intelligent Characterization of Hydration Degree in Cementitious Materials: A Research Plan}
\author{Non-metallic Materials AI+ Laboratory -- Open Fund Project\\ \small China Building Materials Academy (CBMA)}
\date{\today}

\begin{document}

\maketitle

\begin{abstract}
Backscattered-electron (BSE) imaging is the established basis for quantitative microstructural analysis of cementitious materials \citep{scrivener2004}, yet its grayscale contrast is degenerate for chemically similar constituents: supplementary cementitious materials such as metakaolin, slag, fly ash and limestone overlap in gray level with hydration products, so that further refinement of single-modality segmentation architectures yields essentially no gain \citep{lothenbach2011,juenger2015}. This project will overcome this modality information ceiling through multimodal BSE-EDS fusion organized in five parts. A paired dataset spanning four binder systems, four replacement levels and six curing ages will be built, comprising at least 3000 same-field BSE-EDS image pairs and 600 pixel-level annotations cross-checked by QXRD, TGA and EDS point analysis. A two-stage network will couple mutual-information-based cross-modal registration \citep{maes1997} with a fusion segmentation stage in which sparse EDS samples propagate over a chemistry-gradient-constrained superpixel graph \citep{achanta2012} and merge with dense convolutional features \citep{ronneberger2015} through multi-scale attention. A phase-separability criterion with uncertainty-driven active sampling will direct where EDS acquisition is needed, micro-phase descriptors will feed predictors of hydration degree and strength, and integration with the CBMA large model will generate structured characterization reports automatically. Quantitative targets are an IoU improvement of at least 20 percentage points for gray-level-overlapping phases at no more than 5\% EDS coverage, segmentation-derived hydration degree within 5\% of TGA, at least a 50\% reduction in EDS machine time at equal accuracy, and open release of the dataset, fine-tuned CBMA weights and the report-generation module.
\end{abstract}

\medskip
\noindent\textbf{Keywords:} BSE imaging, EDS, multimodal fusion, cementitious materials, phase segmentation, active sampling, hydration degree

\begin{center}\fbox{\parbox{0.9\linewidth}{\small\itshape Research-plan report generated by the paperflow pipeline from the project application. Figures use illustrative synthetic data; all quantitative values are proposed targets, not achieved results.}}\end{center}

\section{Introduction}

Backscattered-electron (BSE) imaging has long served as the workhorse of quantitative microstructure analysis for cementitious materials, because atomic-number contrast on polished sections allows anhydrous clinker, hydration products and porosity to be distinguished and measured \citep{scrivener2004}. At the same time, partial replacement of Portland clinker by supplementary cementitious materials (SCMs) such as slag, fly ash, metakaolin and limestone has become the principal route to low-carbon binders, making blended systems the norm rather than the exception in modern practice \citep{lothenbach2011,juenger2015}. This shift exposes a fundamental limitation of BSE-based analysis: many SCM particles and their reaction rims occupy grayscale ranges that overlap those of the hydrates, so chemically distinct phases become degenerate in a single intensity channel. Deep segmentation networks, from fully convolutional architectures to encoder-decoder designs such as U-Net, have substantially advanced microstructural image analysis \citep{long2015,ronneberger2015}, yet no architecture can recover information that the imaging modality does not encode; for gray-level-overlapping phases, further tuning of BSE-only models yields essentially no gain, revealing an information ceiling intrinsic to the modality itself. Energy-dispersive spectroscopy (EDS) carries precisely the missing chemical discrimination, but dense elemental mapping at segmentation-grade resolution is prohibitively slow for routine or large-area characterization. This tension motivates the central idea of the present project: fusing high-resolution BSE morphology with sparse, intelligently placed EDS measurements through a registration-aware, graph-based fusion network, so that a small fraction of chemical sampling lifts pixel-level phase identification above the BSE ceiling and enables accurate, efficient characterization of hydration degree in blended cementitious systems. This report presents the research plan of the funded project, covering dataset construction, two-stage fusion modeling, uncertainty-driven active sampling, micro-phase-based performance prediction and integration with the CBMA large-model platform.

\section{Methodology}

The experimental programme will produce a systematically designed multimodal dataset covering P.I 42.5 Portland cement blended with slag, fly ash and limestone at 0, 20, 40 and 60 wt\% replacement (w/b = 0.4), sampled at 1, 3, 7, 28, 90 and 180 days to span the hydration trajectories of both clinker and supplementary cementitious phases \citep{lothenbach2011,juenger2015}. For every field of view, high-resolution BSE micrographs will be acquired together with same-field EDS maps of Ca, Si, Al, Fe, Mg and S, yielding at least 3000 paired images, of which at least 600 will carry pixel-level phase annotations cross-checked against QXRD, TGA and EDS point analysis in line with established quantitative BSE practice \citep{scrivener2004}.

The two-stage network first resolves cross-modal registration: rigid alignment will maximize mutual information \citep{maes1997} under joint constraints matching BSE edges to element-gradient edges, after which a learned deformable field corrects stage drift and charging-induced distortion; a field of view advances to segmentation only when residuals at manual landmarks fall below one pixel. In the fusion stage, a dense convolutional encoder processes the BSE image while sparse EDS sampling points are injected as nodes of a superpixel graph whose boundaries are constrained by chemistry gradients \citep{achanta2012}; a graph neural network propagates chemical evidence across superpixels, and multi-scale attention fuses graph and convolutional features. Purely BSE-driven deep baselines \citep{ronneberger2015,long2015} will be trained on identical splits, with a target improvement of at least 20 percentage points in IoU for gray-level-overlapping phases at no more than 5\% EDS coverage. Figure 1 illustrates the pipeline: panel (a) shows a synthetic BSE micrograph, (b) sparse EDS sampling of the same field, (c) the chemistry-constrained superpixel graph with injected nodes, and (d) the resulting phase segmentation.

A phase-separability criterion built on pixel-level uncertainty from Monte-Carlo dropout or deep ensembles will cluster high-entropy regions so that the model designates where further EDS acquisition is most informative, targeting at least a 50\% reduction of instrument time at equal accuracy. Finally, segmentation-derived hydration degree, porosity, phase fractions and phase-adjacency features, independently validated by TGA, QXRD and MIP, will feed predictors of hydration degree and strength, with leave-one-system-out extrapolation testing whether these micro intermediate variables generalize beyond mix-design-only baselines.

\begin{figure}[htbp]
\centering
\includegraphics[width=0.92\linewidth]{figures/fig1_pipeline.png}
\caption{Illustrative overview of the multimodal BSE-EDS fusion pipeline on a synthetic blended-cement microstructure. (a) Synthetic backscattered-electron (BSE) micrograph in which metakaolin and slag occupy nearly identical gray levels, illustrating the gray-level degeneracy that defeats BSE-only segmentation; pores appear dark and clinker bright. (b) Sparse EDS sampling points (coverage of at most 5\% of the field) colored by true phase. (c) SLIC superpixel tessellation of the BSE image with the region-adjacency graph used by the graph branch: superpixels containing an EDS sample appear as filled phase-colored nodes, unsampled superpixels as small gray nodes to be completed by graph propagation. (d) Ground-truth phase map of the same field in the same palette, with thin white phase boundaries. All panels show synthetic, illustrative data generated from gaussian-filtered random fields; they are not measurements.}
\label{fig:pipeline}
\end{figure}

\begin{figure}[htbp]
\centering
\includegraphics[width=0.92\linewidth]{figures/fig2_targets.png}
\caption{Illustrative quantitative targets of the project; the curves and bars are schematic targets, not measured results. (a) Intended IoU of the gray-overlapping phases (metakaolin, slag) versus EDS coverage: the BSE-only baseline stays near 0.55 regardless of architecture tuning, fusion with random sparse sampling improves slowly, and fusion with uncertainty-driven active sampling is targeted to exceed 0.75 before 5\% coverage, meeting the +20 percentage-point target at no more than 5\% coverage. (b) Intended EDS machine time required to reach equal target accuracy, normalized to dense full-field mapping (100\%): active sampling is targeted to require at most 50\% of the dense-mapping time, a reduction of at least 50\%.}
\label{fig:targets}
\end{figure}

\section{Expected Results}

The project defines four quantitative targets, each paired with an explicit validation protocol, and Figure 2 illustrates the two central ones: panel (a) depicts the anticipated intersection-over-union (IoU) as a function of EDS coverage for the BSE-only baseline, a random-sampling fusion variant and the proposed active-sampling fusion network, while panel (b) compares the corresponding EDS machine time. The first target concerns segmentation accuracy: for the gray-level-overlapping phases metakaolin and slag, the fusion model is expected to raise IoU by at least 20 percentage points over a BSE-only baseline implemented with standard fully convolutional and U-Net architectures \citep{long2015,ronneberger2015}, while consuming no more than 5\% EDS pixel coverage; both models will be evaluated on the held-out pixel-annotated test set whose labels are cross-checked by EDS point analysis. The second target concerns physical fidelity: hydration degree computed from the segmented phase maps, following the established convention for quantitative BSE microstructure analysis \citep{scrivener2004}, is required to agree with independent thermogravimetric measurements within 5\% across all mixtures and curing ages. The third target quantifies experimental efficiency: at matched segmentation accuracy, uncertainty-guided active sampling is designed to reduce EDS acquisition time by at least 50\% relative to uniform random sampling, verified by logging actual dwell and mapping times on identical fields of view. The fourth target specifies deliverables: an open dataset of at least 3000 paired BSE-EDS image sets including at least 600 pixel-level annotated images, the fine-tuned CBMA vision backbone weights and the automatic report-generation module will all be deposited in the CBMA inorganic non-metallic materials data node in documented, reusable formats.

\section{Conclusion}

This research plan confronts the grayscale degeneracy that has long constrained backscattered-electron imaging of blended cementitious systems \citep{scrivener2004}, in which supplementary cementitious materials remain difficult to distinguish from hydration products by morphology alone \citep{lothenbach2011}. By coupling sub-pixel BSE-EDS registration with graph-convolutional fusion segmentation, and by letting pixel-level uncertainty designate where elemental acquisition is most informative, the plan is designed to raise the IoU of gray-level-overlapping phases by at least twenty percentage points at no more than five percent EDS coverage, while halving instrument time at equal accuracy. Conceptually, the framework recasts EDS from a passive verification instrument into an actively directed supply of chemical information, so that the model itself decides where the experiment should look next. The planned deliverables, comprising an open dataset of paired and annotated micrographs, fine-tuned CBMA vision weights, and an automatic report-generation module, will constitute a reusable intelligent-characterization capability serving both the laboratory and the CBMA inorganic non-metallic materials data node. Future work will extend the framework toward cross-system generalization and fully automated microstructural reporting, moving quantitative SEM characterization from expert-bound practice toward a routinely deployable service.


\bibliographystyle{plainnat}
\bibliography{references}

\end{document}
